% =============================================================================
% Section 7: Obstruction analysis
% Why standard analytic upgrades do not turn the barrier into closure.
% =============================================================================

\section{Why standard analytic upgrades do not close the frontier}
\label{sec:obstructions}

The audit chain of \cref{sec:csp-audit} establishes the
finite-local insufficiency. A natural next attempt is to add an
analytic upgrade --- bound the count of survivors via the standard
sieve machinery (Selberg, Shiu, Nair--Tenenbaum, GPY, Maynard,
Tao--Ter\"av\"ainen) and combine with finite verification to close
the frontier. This section explains, with explicit lemmas where
appropriate, why the standard analytic upgrades are also blocked.

The four convergent obstructions are:

\begin{enumerate}[leftmargin=*, label=(\textbf{O\arabic*})]
\item the \emph{Selberg upper-bound divergence} for fixed-$K$
  multidimensional sieves (\cref{sec:selberg-divergence});
\item the \emph{$L_1$ vs $L_\infty$ gap} between average and pointwise
  divisor estimates (\cref{sec:L1-Linf});
\item the \emph{$k \le 2$ correlation ceiling} for multi-point
  divisor correlations (\cref{sec:correlation-ceiling});
\item the \emph{wrong polarity} of prime-detection sieves applied to
  divisor-suppression problems (\cref{sec:wrong-polarity}).
\end{enumerate}

Each obstruction is independently sufficient to block closure of
\cref{prob:e647} via the corresponding family of methods, when the
target is a fixed finite window or an almost-all/density upgrade. We
address them in order. The corrected Stage--1 framing leaves one
separate theorem-design path open: a pointwise shifted-prime tail
theorem on the $41 \cdot 529$ explicit open sub-progressions. This
Shape--C path is recorded below as an open target, not as a closure
result and not as a consequence of the present literature.

\subsection{The Selberg upper-bound divergence}
\label{sec:selberg-divergence}

The natural analytic attempt is to sieve the \PLEclass-survivor count
on a residue class with a multidimensional Selberg $\lambda^2$ weight
of the form developed by Maynard \cite{Maynard2013} and refined by
Polymath~8b \cite{Polymath8b2014}. Under the standard Selberg
machinery, one obtains an upper bound of the form
\[
\#\bigl\{ Y \le X : Y \equiv u \pmod{323}, \;
\tauf(360360 Y - k) \le k + 2 \;\forall\, 1 \le k \le K \bigr\}
\;\le\; C \cdot \frac{X}{(\log X)^K}
\]
for an explicit constant $C = C(u, K)$ depending on the singular
series and the choice of sieve weight.

\begin{lemma}[Selberg upper-bound divergence]
\label{lem:selberg-divergence}
Fix $K \ge 1$ and $C > 0$. The function
\[
f(x) \;=\; C \cdot \frac{x}{(\log x)^K}, \qquad x > 1,
\]
satisfies
\[
f'(x) \;=\; C \cdot \frac{\log x - K}{(\log x)^{K+1}} ,
\]
and is therefore strictly decreasing on $1 < x < e^K$, with global
minimum
\[
f(e^K) \;=\; C \cdot e^K \cdot K^{-K} \;=\; C \cdot \biggl( \frac{e}{K} \biggr)^{K},
\]
and strictly increasing for $x > e^K$. In particular,
\[
\lim_{x \to \infty} f(x) \;=\; +\infty .
\]
\end{lemma}

\begin{proof}
Direct computation. The derivative vanishes at $\log x = K$, and
$f(x) / x = C / (\log x)^K \to 0$ slowly as $x \to \infty$, but
$f(x) \to \infty$ since $x$ grows super-polynomially in $\log x$.
\qed
\end{proof}

\begin{corollary}[Sparsity bounds do not close]
\label{cor:sparsity-not-closure}
Let $S \subset \N$ be a set such that, for some $C > 0$ and $K \ge 1$,
\[
\#\bigl( S \cap [1, X] \bigr) \;\le\; C \cdot \frac{X}{(\log X)^K}
\quad \text{for all } X \ge X_0 .
\]
Then $\#(S \cap [X, 2X]) \to \infty$ as $X \to \infty$. In
particular, no such bound certifies the eventual emptiness of $S$.
\end{corollary}

\begin{proof}
$\#(S \cap [X, 2X]) \le f(2X) - f(X) + O(1)$ where $f$ is as in
\cref{lem:selberg-divergence}. For $X > e^K$, $f$ is strictly
increasing and grows like $X / (\log X)^K \to \infty$, so the
dyadic count diverges with $X$. Closure of $S$ would require
$\#(S \cap [X, 2X]) = 0$ for all $X \ge X_0$, which is incompatible
with $\#(S \cap [X, 2X]) \to \infty$.
\qed
\end{proof}

\Cref{cor:sparsity-not-closure} formalizes the slogan
\emph{sparsity is not emptiness}. A bound of the form
$\#(S \cap [1, X]) \le X^{1 - \delta}$, or
$\#(S \cap [1, X]) \le X / (\log X)^A$ for any fixed $A$, allows
$S$ to grow without bound. To certify closure, one needs either a
deterministic finite cover or a dyadic-interval upper bound that is
strictly less than $1$ for all sufficiently large $X$.

\subsection{Why density-zero is not emptiness: the \texorpdfstring{$L_1$ vs $L_\infty$}{L1 vs Linfty} gap}
\label{sec:L1-Linf}

A second natural attempt is to invoke average estimates for the
divisor function on the residual frontier in the form of Shiu
\cite{Shiu1980}, Nair--Tenenbaum \cite{NairTenenbaum1998}, or
Henriot \cite{Henriot2012}. Shiu's theorem gives, for a multiplicative
function $f$ in an arithmetic progression and a short interval,
\[
\sum_{\substack{x < n \le x + y \\ n \equiv a \pmod q}} f(n)
\;\ll\; \frac{y}{\phi(q) \log x}
\exp \biggl( \sum_{\substack{p \le x \\ p \nmid q}} \frac{f(p)}{p} \biggr) .
\]
For $f = \tauf$ this yields the average $\sum_{n \in [x, x+y], n \equiv a \pmod q} \tauf(n) \asymp y \log x / \phi(q)$, with an analogous lower bound when the conditions of Shiu's theorem are met.

These results control the \emph{$L_1$ average} of $\tauf$ over the
sequence. They do not control the pointwise minimum.

\begin{remark}[Pointwise barrier failure requires $L_\infty$ data]
The $\tauf$-barrier inequality demands
$\tauf(n - k) \le k + 2$ \emph{pointwise for every $1 \le k \le K$}.
A lower bound on $\sum_{m \in [x, x+y]} \tauf(m)$ does not imply
that any individual $\tauf(m)$ is large, since a small number of
highly divisor-rich integers can sustain the average while the
remaining integers have low $\tauf$. To rule out a survivor
candidate $Y \equiv u \pmod{323}$, one needs an $L_\infty$
estimate: that some $k$ in the window forces $\tauf(360360 Y - k) > k+2$.
The $L_1$-average tools do not produce this, even under strong
conditional hypotheses.
\end{remark}

This is the second obstruction, and the no-closure-via-correlation
conclusion is the convergent verdict across the four obstruction
families enumerated in this section.

\subsection{The $k \le 2$ correlation ceiling}
\label{sec:correlation-ceiling}

The most ambitious analytic attempt would be to control multi-point
$\tauf$-correlations, e.g.\
\[
\sum_{n \le X} \prod_{i = 1}^{r} \tauf(n - k_i) \cdot
\mathbf{1}\bigl[ n \equiv u \pmod{323} \bigr],
\]
and rule out the simultaneous occurrence of low $\tauf$-values across
the protected window. This is the natural setting of the
contemporary correlation-of-arithmetic-functions literature
\cite{TaoTeravainen2025, MRT2019, MSTT2024_I, MSTT2024_II}.

\begin{remark}[The $k \le 2$ correlation ceiling]
\label{rem:correlation-ceiling}
Pointwise unconditional asymptotic formulas for two-point
correlations of multiplicative functions exist in special cases
(Estermann--Ingham for $\tauf \cdot \tauf$ and analogues). Three-point
correlations of $d_k$'s are open: Matom\"aki--Radziwi\l{}\l --Tao
\cite{MRT2019} explicitly state $\sum d_3(n) d_3(n + h)$ is ``a
notorious open problem,'' unresolved even for $o(X)$ error. The
$\tauf$-barrier of \cref{prob:e647} requires simultaneous control of
$\tauf$-values across a window of length $K \ge 3$, which is a
multi-point conjunctive condition --- exactly where the literature
ends.

The ``almost-all'' framework of MSTT-I/II
\cite{MSTT2024_I, MSTT2024_II} controls $f - f^\sharp$ on average
over short intervals, with an exceptional set of measure
$X (\log X)^{-A}$. This exceptional set is vastly larger than any
single residue class modulo $M = 46189$; an almost-all statement
cannot be promoted to pointwise without a fundamentally new
ingredient. Per the analysis in \cite[footnote 4]{MRT2019}: even
under the Generalized Lindel\"of Hypothesis, GRH, or
Elliott--Halberstam, the methods cannot work for short-interval
length $H$ slower than $\log X$, because the $h = 0$ term dominates
the averages.
\end{remark}

\subsection{Wrong polarity}
\label{sec:wrong-polarity}

A fourth obstruction concerns the polarity of the available sieve
machinery. Maynard-style multidimensional sieves
\cite{Maynard2013, Maynard2014} are tuned to \emph{detect} primes
--- they extract integers $n$ such that many of $n + h_1, \ldots,
n + h_K$ are simultaneously prime, with an explicit lower bound on
the count of such $n$. The $\tauf$-barrier requires the opposite:
\emph{suppress} divisors, ruling out integers $n$ such that any
$\tauf(n - k) > k + 2$.

\begin{remark}[Wrong polarity]
\label{rem:wrong-polarity}
The natural inversions of prime-detection sieves yield results in
the opposite direction --- counts of $n$ with \emph{many} prime
factors (or \emph{many} divisors) on shifts, not bounds on $n$ with
\emph{few}. Maynard 2014's Theorem 3.3 produces many $n$ such that
$n + h_1, \ldots, n + h_{m+1}$ are simultaneously prime in a fixed
arithmetic progression --- exactly the opposite direction from a
$\tauf$-suppression theorem. The Tao--Ter\"av\"ainen 2025 result on
problem \#248 \cite{TaoTeravainen2025} establishes
$\Omega(n + k) \le C k$ on infinitely many $n$ --- a multiplicative
$Ck$ budget rather than the additive $k + 2$ budget required by
\cref{prob:e647}.
Tao--Ter\"av\"ainen explicitly note that the corresponding
$\tauf(n-k) \le k+2$ barrier-case variant is beyond their methods;
Pilatte's logarithmic two-point Chowla saving \cite{Pilatte2024}
is one of the inputs that makes their two-point correlation machinery
effective, but it is still an averaged two-point statement rather
than a prime-conditioned pointwise divisor lower bound.

The structural reason is the Maynard variance ratio
$M_k \approx \log k - 2 \log \log k$, forced by the Mertens estimate
$\sum_{p \le x} 1/p = \log \log x + O(1)$. To achieve the tighter
additive $k + 2$ budget would require sub-Mertens variance bounds,
which are not available within standard sieve theory.
\end{remark}

\subsection{Synthesis: closure requires emptiness, not sparsity}
\label{sec:synthesis-no-closure}

The four obstructions of this section converge on a single point.
Closure of \cref{prob:e647} on the residual frontier requires
either:

\begin{enumerate}[leftmargin=*, label=(\roman*)]
\item a \emph{deterministic finite cover}: a finite list of
  congruence patterns whose union forces $\tauf(n - k) > k + 2$ for
  some $k$ at every $n$ in the residue class above an explicit
  threshold; or
\item a \emph{dyadic-interval upper bound strictly less than $1$}:
  $\#(S \cap [X, 2X]) < 1$ for all $X \ge X_0$, where $S$ is the
  survivor set.
\end{enumerate}

The audit chain of \cref{sec:csp-audit} establishes that the
finite-local primitives in our toolbox cannot produce the first.
The Selberg divergence (\cref{lem:selberg-divergence}), the
$L_1$-vs-$L_\infty$ gap (\cref{sec:L1-Linf}), the $k \le 2$
correlation ceiling (\cref{rem:correlation-ceiling}), and the
wrong-polarity issue (\cref{rem:wrong-polarity}) collectively
establish that the contemporary analytic-NT tools cannot produce
the second.

\begin{remark}[What remains genuinely open]
\label{rem:what-remains-open}
Four classes of methods are not addressed by the obstructions of
this section:
\begin{itemize}[leftmargin=*]
\item Erd\H{o}s covering systems with non-prime-power moduli
  \cite{Hough2015, BBMST2022}, when the cover is genuinely
  combinatorial rather than CRT-decomposable into prime-residual
  kernels;
\item methods deriving from a future pointwise three-or-more-point
  divisor correlation theorem at $q = M$ (equivalently, a positive
  resolution of the open Conrey--Gonek--Keating asymptotic at
  $q = M$);
\item a pointwise shifted-prime interval theorem on the explicit
  Stage--1 open sub-progressions. In the notation of
  \cref{sec:structural-reduction}, write
  \[
    P = F_1(r,s,u) = A u + B(r,s)-1,
    \qquad A = 2520 \cdot 46189 \cdot 529 = 61573632120.
  \]
  A sufficient Shape--C target would be: for some
  $\varepsilon > 1/\log 2$ and $C>0$, every sufficiently large prime
  $P=F_1(r,s,u)$ with $r$ in the 41 open residues and
  $0 \le s < 529$ has a shift
  \[
    1 \le d \le C \log P \sqrt{\log\log P}
  \]
  such that $\omega(P-d) \ge \varepsilon \log\log P$.
  Then $\tauf(P-d) \ge 2^{\omega(P-d)}
       \ge (\log P)^{\varepsilon\log 2}$, which eventually exceeds
  the linear budget $d+3$ after translating back to $k=d+1$.
  Lau's conjectural interval framework \cite{Lau2026} points at this
  scale, but the restricted prime-endpoint theorem above is not
  proved by Lau, Tao--Ter\"av\"ainen, or Pilatte;
\item methods that establish closure via a strictly-greater-than-zero
  density of \emph{counterexamples}, by exhibiting one explicitly ---
  a path that would refute Erd\H{o}s' conjecture but is consistent
  with the structural reduction of \cref{thm:stage1-reduction}.
\end{itemize}
The first three are recognized as open problems in their own right;
the last is the status quo of the search. The shifted-prime interval
target is the only route above that directly matches the current
Lean-facing non-uniform wrapper: it is pointwise in the prime channel,
may use a growing $K(r,s,u)$, and is restricted to the 21689 explicit
sub-progressions rather than all primes. It should therefore be
treated as a partial-open Shape--C theorem-design target, not as an
exception to the density/average obstructions of this section.
\end{remark}

\bigskip

The combination of \cref{sec:csp-audit,sec:obstructions} constitutes
the barrier theorem in the title: the natural finite-local
forced-divisor methods do not close the residual frontier, and the
natural analytic-NT upgrades available with current technology do
not close the frontier either. Erd\H{o}s' 1979 statement that the
problem is ``hopeless with our present methods'' is, on the
evidence of this audit, accurate and structural.

% =============================================================================
% End of section 7
% =============================================================================
